\section{Immutable Parent Pointers}

\subsection{Formal Definition}

The \textbf{Recursive Spawning} framework establishes an immutable directed graph structure between agents and their creators. At its core lies the \textbf{ParentPointer} relation, which records the unique identity of each agent's originating parent at the moment of spawn.

\begin{definition}[ParentPointer Relation]
Let $p$ denote a parent agent and $c$ denote a child agent. The relation
\[
\text{ParentPointer}(p, c)
\]
is established at the moment of spawn and \textbf{never modified} thereafter. The tuple $(p, c)$ is stored as an immutable record in the agent's genesis block.
\end{definition}

\begin{axiom}[Immutability Axiom]
Once $\text{ParentPointer}(p, c)$ is recorded, the relation $(p, c)$ cannot be altered, replaced, or deleted by any subsequent operation, including:
\begin{itemize}
    \item Administrative action
    \item Runtime reconfiguration
    \item Agent self-modification
    \item External system override
\end{itemize}
The only exception is agent termination, which preserves the record but marks the agent as \texttt{DEACTIVATED}.
\end{axiom}

\subsection{Chain Construction}

The complete ancestry of any agent forms a recursively defined chain terminating at the root human. We define the chain construction recursively:

\begin{definition}[Chain Function]
Define $\text{Chain}(a)$ for any agent $a$ recursively:
\[
\text{Chain}(a) \equiv \text{ParentPointer}(\text{parent}(a), a) \;\cup\; \text{Chain}(\text{parent}(a))
\]
where $\cup$ denotes set union of pointer records. The empty set $\emptyset$ appears as the base case for the root human.
\end{definition}

\begin{definition}[Base Case: Root Human]
The root human entity---denoted $h_\emptyset$---has no parent. Formally:
\[
\text{parent}(h_\emptyset) = \perp \quad \text{(null/void)}
\]
Consequently:
\[
\text{Chain}(h_\emptyset) = \emptyset
\]
This establishes the terminal boundary of any agent ancestry chain.
\end{definition}

\begin{example}[Chain Traversal]
Consider an agent $a_3$ spawned by agent $a_2$, which was itself spawned by agent $a_1$, which was spawned by the root human $h_\emptyset$:
\[
\text{Chain}(a_3) = \{
    \text{ParentPointer}(a_2, a_3),
    \text{ParentPointer}(a_1, a_2),
    \text{ParentPointer}(h_\emptyset, a_1)
\}
\]
This chain contains exactly three immutable records, reflecting the complete spawn lineage.
\end{example}

\subsection{Structural Properties}

The parent pointer structure exhibits several critical properties that enforce the integrity of the recursive spawning model.

\begin{property}[Acyclicity]
The parent pointer relation guarantees a directed acyclic graph (DAG) structure. No agent can appear as its own ancestor. Formally:
\[
\forall a: a \notin \text{Chain}(a)
\]
This prevents circular spawn relationships and maintains a strict temporal ordering.
\end{property}

\begin{property}[Uniqueness]
Each agent has exactly one parent pointer. While spawn events may occur in parallel, each child $c$ receives precisely one assignment $\text{ParentPointer}(p, c)$ at creation. Formally:
\[
\forall c: |\text{ParentPointer}(p, c)| \leq 1
\]
If $c$ is not the root human, $|\text{ParentPointer}(p, c)| = 1$.
\end{property}

\begin{property}[Temporal Ordering]
Parent pointers preserve causal ordering. Since parent assignment occurs atomically at spawn, we guarantee:
\[
\text{spawn\_time}(p) < \text{spawn\_time}(c) \quad \forall\; \text{ParentPointer}(p, c)
\]
\end{property}

\begin{figure}[htbp]
\centering
\inputcedence{}{Chain Structure in Recursive Spawning}
\caption{Illustration of the Chain structure showing immutable parent pointers and the recursive composition of ancestry records. Each node represents an agent; directed edges represent immutable parent pointer relations. The chain terminates at the root human with null parent.}
\label{fig:chain-structure}
\end{figure}

\subsection{Immutability Enforcement: Fact Layer}

The Fact Layer serves as the authoritative enforcement mechanism for parent pointer immutability. All modifications to agent state flow through the Fact Layer, which implements strict validation and logging protocols.

\begin{definition}[Fact Layer Role]
The Fact Layer is the sole authoritative record-keeper for parent pointer assignments. It provides:
\begin{enumerate}
    \item \textbf{Write-once registration} of parent pointers at spawn time
    \item \textbf{Read-only access} to parent pointer records post-creation
    \item \textbf{Modification blocking} for any attempt to alter existing pointer relations
    \item \textbf{Audit logging} of all pointer-related operations
\end{enumerate}
\end{definition}

\begin{invariant}[Fact Layer Immutability Lock]
Any attempt to execute $\text{ParentPointer}(p', c)$ where $c$ already has a recorded parent results in:
\[
\text{REJECT} \quad \text{if} \quad \exists\; p: \text{ParentPointer}(p, c) \;\text{is already registered}
\]
The Fact Layer refuses to overwrite, update, or delete existing parent pointers.
\end{invariant}

\begin{algorithm}[Parent Pointer Registration]
\begin{enumerate}
    \item Spawn event triggers parent pointer creation
    \item Fact Layer receives $\text{ParentPointer}(p, c)$ registration request
    \item Fact Layer queries existing registry for any prior parent of $c$
    \item \textbf{If prior parent exists}: return \texttt{ERROR\_DUPLICATE\_CHILD} and abort spawn
    \item \textbf{If no prior parent}: record $(p, c)$ as immutable tuple, timestamp with $\tau_\text{spawn}$
    \item Return $\texttt{REGISTERED}_p$ confirmation to Purpose Layer
    \item Chain entry appended to genesis block; genesis block sealed
\end{enumerate}
\end{algorithm}

\begin{property}[Genesis Block Sealing]
Upon successful registration, the genesis block containing the parent pointer record is cryptographically sealed. The seal ensures:
\begin{itemize}
    \item Tamper-evident storage via hash chaining
    \item Non-repudiation of spawn event
    \item Preservation of pointer record across agent migrations
\end{itemize}
\end{property}

The Fact Layer blocks modification through architecture-level enforcement:

\begin{itemize}
    \item \textbf{No UPDATE API} for parent pointer fields
    \item \textbf{No DELETE API} for parent pointer records
    \item \textbf{No OVERWRITE semantics} in pointer storage
    \item \textbf{Auditing} of all parent pointer access attempts (authorized or not)
\end{itemize}

\subsection{Spawn Authorization: Purpose Layer}

The Purpose Layer initiates spawn events and requests parent pointer assignment from the Fact Layer. It carries the burden of proving valid authorization before any parent pointer is recorded.

\begin{definition}[Purpose Layer Role in Spawn]
The Purpose Layer:
\begin{enumerate}
    \item Receives or generates spawn authorization credentials
    \item Validates that the proposed spawn adheres to agent purpose and policy
    \item Submits parent pointer registration request to Fact Layer
    \item Receives confirmation of immutable parent pointer registration
    \item Proceeds with agent instantiation only upon confirmed registration
\end{enumerate}
\end{definition}

\begin{protocol}[Spawn Authorization Protocol]
\begin{enumerate}
    \item Spawn request arrives at Purpose Layer with target agent template
    \item Purpose Layer evaluates spawn eligibility against policy rules:
    \begin{itemize}
        \item Purpose alignment check
        \item Resource allocation validation
        \item Authorization scope verification
        \item Cloning vs.\ de novo spawn determination
    \end{itemize}
    \item If eligible, Purpose Layer constructs registration payload:
    \begin{itemize}
        \item $\text{parent}(c)$ identifier
        \item $c$ proposed identifier
        \item Spawn authorization evidence (token/signature)
        \item Timestamp $\tau_\text{auth}$
    \end{itemize}
    \item Purpose Layer submits $\text{ParentPointer}(\text{parent}(c), c)$ to Fact Layer
    \item Fact Layer verifies no prior parent assignment exists for $c$
    \item On success: parent pointer recorded, genesis block sealed
    \item On failure: spawn aborts, parent pointer not created
\end{enumerate}
\end{protocol}

\begin{invariant}[Authorization Required]
No agent may be instantiated without a corresponding parent pointer registration. The parent pointer registration is a prerequisite, not a post-condition, of spawn completion. Formally:
\[
\text{SpawnComplete}(c) \implies \exists\; p: \text{ParentPointer}(p, c) \;\text{is registered}
\]
\end{invariant}

The Purpose Layer cannot bypass Fact Layer immutability guarantees. Even with administrative override privileges, the Fact Layer architecture rejects any attempt to modify an existing parent pointer. Override authority does not extend to altering recorded history.

\subsection{Chain Traversal Algorithm}

Given an agent identifier $a$, we can traverse its ancestry chain recursively to reconstruct its spawn lineage. The traversal algorithm terminates at the root human where the parent field is null.

\begin{algorithm}[Chain Traversal]
\begin{verbatim}
FUNCTION ChainTraversal(agent a):
    result ← empty list
    current ← a
    
    WHILE parent(current) ≠ ⊥:
        pp ← ParentPointer(parent(current), current)
        result.APPEND(pp)
        current ← parent(current)
    
    RETURN result
\end{verbatim}
\end{algorithm}

\begin{algorithm}[Recursive Chain Construction]
\begin{verbatim}
FUNCTION Chain(a):
    IF parent(a) = ⊥:
        RETURN ∅
    ELSE:
        RETURN ParentPointer(parent(a), a) ∪ Chain(parent(a))
\end{verbatim}
\end{algorithm}

Both implementations produce identical results. The recursive version more directly mirrors the formal definition; the iterative version avoids stack overflow for deep chains.

\begin{complexity}[Chain Traversal]
Let $n = |\text{Chain}(a)|$ represent the number of ancestry levels. Then:
\begin{itemize}
    \item \textbf{Time complexity}: $\mathcal{O}(n)$
    \item \textbf{Space complexity}: $\mathcal{O}(n)$ for storing all pointer records
    \item \textbf{No duplication}: Each pointer record appears exactly once in the result
\end{itemize}
\end{complexity}

\begin{note}
Chain traversal is a read-only operation. It accesses immutable records without modifying any state. Concurrent traversal operations do not interfere with one another.
\end{note}

\begin{example}[Traversal Walkthrough]
Given agent $a_3$ with parents $a_2 \to a_1 \to h_\emptyset$:
\begin{enumerate}
    \item $\text{ChainTraversal}(a_3)$ begins with \texttt{current} = $a_3$
    \item $\text{parent}(a_3) = a_2 \neq \perp$; append $\text{ParentPointer}(a_2, a_3)$
    \item $\text{current} \leftarrow a_2$
    \item $\text{parent}(a_2) = a_1 \neq \perp$; append $\text{ParentPointer}(a_1, a_2)$
    \item $\text{current} \leftarrow a_1$
    \item $\text{parent}(a_1) = h_\emptyset \neq \perp$; append $\text{ParentPointer}(h_\emptyset, a_1)$
    \item $\text{current} \leftarrow h_\emptyset$
    \item $\text{parent}(h_\emptyset) = \perp$; terminate
    \item Return list of three pointer records
\end{enumerate}
\end{example}

\subsection{Policy Implications of Immutability}

The immutability of parent pointers carries significant implications for system policy and governance.

\begin{policy}[Chain of Custody]
Because parent pointers cannot be altered post-spawn, the complete ancestry chain serves as an unbroken chain of custody for accountability. Any agent's lineage can be audited unambiguously by traversing $\text{Chain}(a)$. This supports:
\begin{itemize}
    \item \textbf{Compliance reporting}: Tracing authorization origin for regulatory requirements
    \item \textbf{Fault analysis}: Determining root cause by examining spawn chains
    \item \textbf{Delegation chains}: Understanding who authorized what in multi-tier spawning
\end{itemize}
\end{policy}

\begin{policy}[No Orphaning]
Agents cannot be disconnected from their ancestry. The immutability guarantee ensures every non-root agent has a recorded parent, preventing orphaned agent records in the system.
\end{policy}

\begin{policy}[Auditability]
Immutable parent pointers enable complete history reconstruction. Because the Fact Layer never permits modification, the recorded spawn history reflects the actual spawn events exactly as they occurred, without post-hoc alterations or retroactive repairs.
\end{policy}

\subsection{Formal Theorems}

We state two fundamental theorems deriving from the immutable parent pointer architecture.

\begin theorem}[Termination Guarantee]
Every chain $\text{Chain}(a)$ constructed via the recursive definition terminates at the root human $h_\emptyset$ in finite depth.
\end theorem}
\begin-proof}
Consider any agent $a$. By the definition of spawn, agents are created at discrete spawn events with strictly increasing timestamps. Since spawn events cannot cycle (acyclicity property), there are finitely many ancestors. Repeatedly applying the parent function must eventually reach an agent with no parent, namely $h_\emptyset$, because the set of all agents is finite and the parent relation creates no cycles. Therefore $\text{Chain}(a)$ terminates in finite depth.
\end-proof}

\begin theorem}[Ancestry Uniqueness]
For any two distinct agents $a$ and $b$ with $a \neq b$, we have $\text{Chain}(a) \neq \text{Chain}(b)$.
\end theorem}
\begin-proof}
Each agent's chain terminates at a unique leaf (itself) combined with its distinct ancestry. Two agents that diverge from a common ancestor share a partial chain but differ at the point where one path includes a child not present in the other's ancestry. Since each agent's own parent pointer is unique to it, the sets of pointer records differ. Therefore the complete chains are not identical.
\end-proof}

These theorems establish that the immutable parent pointer structure yields a well-defined, unambiguous ancestry graph with guaranteed termination and unique identification of each agent's lineage.